\begin{center}{\large\textbf{The Greenhouse Effect}}\end{center}

\textit{In 2021, Syukuro Manabe and Klaus Hasselmann were awarded half of the
Nobel Prize in Physics for their work in modeling Earth's climate and
accurately predicting the global warming caused by human industrial
activities. In this problem, we will examine a simple model of global warming
due to the greenhouse effect. The greenhouse gases alter the optical properties
of the Earth's atmosphere in transmitting or absorbing Earth's infrared
radiation, resulting in a rise in the average temperature of the planet.}

All objects, at different temperatures, emit thermal radiation. The quantity
$u(\lambda, T)d\lambda$ indicates the thermal radiative power per unit area of
an object at temperature $T$ between the wavelengths $\lambda$ and
$\lambda + d\lambda$. According to Planck’s theory of blackbody radiation, we
have
\begin{equation}
u(\lambda, T) = \frac{2\pi hc^2}{\lambda^5}\,\frac{1}{\exp(\,\frac{hc}{\lambda k_\mathrm{B}T}\,)\ -1}\,, \tag{1}
\end{equation}
in which $hc = 1.\,24 \times 10^3\,\mathrm{eV\cdot nm}$ and
$k_\mathrm{B} = 8.\,62 \times 10^{-5}\ \mathrm{eV\,/K}$. The wavelength
corresponding to the maximum of $u(\lambda, T)$ comes from the relation
$\lambda_\mathrm{max}T = b$ (Wien’s displacement law). Indeed, using equation
(1), it can be shown that $b = \frac{hc}{x_\mathrm{m}k_\mathrm{B}}$, where the
dimensionless quantity $x_\mathrm{m}$ is the non-trivial root of an equation of
the form $f(x) = 0$; you are asked to find the function $f(x)$ in one of the
following tasks. Total radiative power per unit area of a blackbody in all
wavelengths is given by the Stephan-Boltzmann law as $U(T) = \sigma T^4$ where
$\sigma = 5.\,67 \times 10^{-8}\,\mathrm{W/m^2K^4}$. Moreover, according to
Kirchhoff’s law of radiation, at thermal equilibrium a body absorbing a certain
fraction of the incident radiation at a specific wavelength, will radiate the
same fraction of the blackbody radiation at that same wavelength.

Throughout this problem assume that the Sun is a blackbody at its average
surface temperature of $T_\mathrm{S} = 5.\,77 \times 10^3\,\mathrm{K}$. The
Sun’s radius is $R_\mathrm{S} = 6.\,96 \times 10^8\,\mathrm{m}$ and the average
distance between the Earth and the Sun is
$d = 1.\,50 \times 10^{11}\,\mathrm{m}$. We denote by
$\tilde{u}_\mathrm{S}(\lambda)$, the spectral solar power radiated into a unit
area of the Earth normal to the direction of radiation. The integral of this
quantity over all wavelengths, i.e.
$S_0 = \int\,\tilde{u}_\mathrm{S}(\lambda)d\lambda$, is called the solar
constant.

In this problem assume that the Earth is in thermal equilibrium and has the
same temperature at all points on its surface. In all parts of the problem,
express the desired quantity in parametric form in terms of the data given in
the problem and then find its numerical value accurate to three significant
figures. The required units are indicated on the answer sheet.

\textbf{A. Earth as a Blackbody}

In this part, consider the Earth’s surface as a blackbody and neglect the
Earth's atmosphere.

A-1\quad Find the solar constant, $S_0$. \hfill 0.6 pt

A-2\quad Find the Earth’s temperature, $T_\mathrm{E}$ . \hfill 0.6 pt

A-3\quad Find the function $f(x)$. \hfill 0.4 pt

A-4\quad Calculate the numerical value of $x_\mathrm{m}$ , and from this value
$x_\mathrm{m}$ , find the value of $b$. \hfill 0.4 pt

A-5\quad Find $\lambda_\mathrm{max}$ for the Sun and the Earth. \hfill 0.2 pt

In figure 1 the functions $\gamma\tilde{u}_\mathrm{S}(\lambda)$ and
$u(\lambda, T_\mathrm{E})$ are plotted versus $\lambda$, where $\gamma$ is a
dimensionless coefficient to rescale $\tilde{u}_\mathrm{S}(\lambda)$ such that
the values of the two peaks coincide.

A-6\quad Determine $\gamma$. \hfill 0.8 pt

\begin{center}\includegraphics[width=0.75\linewidth]{figures/IPhO_2024_Q1_fig1.png}\end{center}

\begin{center}\textit{Figure 1 - The plot of $u(\lambda, T_\mathrm{E})$ (red)
and $\gamma\tilde{u}_\mathrm{S}(\lambda)$ (blue) versus $\lambda$}\end{center}

\textbf{B. The Greenhouse Effect}

In this part, we introduce a simple model in which the Earth’s atmosphere is
modeled as a thin layer at a small distance above the Earth’s surface so that
the difference between the area of the atmosphere’s layer and the area of the
Earth's surface can be neglected (see figure 2). In what follows assume that
the major part of the thermal radiation from the Earth and the Sun are emitted
at wavelengths near the $\lambda_\mathrm{max}$ for each one. Also assume that
the “atmosphere layer” reflects a fraction $r_\mathrm{A} = 0.\,255$ of the
\textbf{visible-ultraviolet} radiation incident from above or below, and
completely transmits the rest. Assume that the atmosphere does not reflect any
part of the \textbf{infrared} radiation, however, it absorbs a fraction
$\varepsilon$ of the \textbf{infrared} radiation and transmits the rest. This
behavior, known as the greenhouse effect, changes the average temperature of
the Earth. The Earth’s surface, on the other hand, reflects a fraction
$r_\mathrm{E}$ of the \textbf{visible-ultraviolet} radiation and absorbs the
rest of this radiation and all the \textbf{infrared} radiation.

\begin{center}\includegraphics[width=0.8\linewidth]{figures/IPhO_2024_Q1_fig2.png}\end{center}

\begin{center}\textit{Figure 2 - Thermal flows between the Earth and the
atmosphere}\end{center}

B-1\quad Assume that $\varepsilon = 1$ and $r_\mathrm{E} = 0$, and calculate
the Earth’s temperature $T_\mathrm{E}$ and the atmosphere’s temperature
$T_\mathrm{A}$. \hfill 1.0 pt

Now assume that $r_\mathrm{E} \neq 0$. In this case, the combined system of
“Earth + atmosphere” reflects a different fraction of the solar radiation,
called “albedo” and denoted by $\alpha$.

B-2\quad Determine the albedo, $\alpha$, in terms of $r_\mathrm{E}$ and
$r_\mathrm{A}$. Then calculate its numerical value assuming
$r_\mathrm{E} = 0.\,102$ (and $r_\mathrm{A} = 0.\,255$). \hfill 1.6 pt

B-3\quad a) Express the Earth’s temperature in terms of $\sigma$, $\alpha$,
$S_0$, and $\varepsilon$.

b) Using the given data and the calculated albedo, find the numerical value of
$\varepsilon$ which leads to the current average temperature of
$T_\mathrm{E} = 288\,\mathrm{K}$ for the Earth. \hfill 1.0 pt

B-4\quad Find $\dfrac{dT_\mathrm{E}}{d\varepsilon}$ and determine by how much
the Earth’s temperature increases if $\varepsilon$ increases by one percent.
\hfill 0.8 pt

Assume $T_\mathrm{A} = 245\ \mathrm{K}$ and $T_\mathrm{E} = 288\ \mathrm{K}$.
These values come from real data and may differ from the results which you have
obtained in the previous tasks. Now suppose that a non-radiative (e. g.
convective) thermal flow $J_\mathrm{NR} = k(T_\mathrm{E} - T_\mathrm{A})$ is
maintained from the Earth to the atmosphere, where $k$ is a constant. The
quantity, $J_\mathrm{NR}$, is the transmitted power per unit area.

B-5\quad Calculate $\varepsilon$ and $k$ in terms of $T_\mathrm{E}$,
$T_\mathrm{A}$, $\sigma$, $\alpha$, and $S_0$. \hfill 1.6 pt

B-6\quad a) Differentiating the equations obtained in part B-5 with respect to
$\varepsilon$, find the two algebraic equations satisfied by
$\frac{dT_\mathrm{A}}{d\varepsilon}$ and $\frac{dT_\mathrm{E}}{d\varepsilon}$.

b) Use these equations to find the numerical value of change in the Earth’s
temperature as a result of a one percent increase in the value of
$\varepsilon$. \hfill 1.0 pt
